\documentclass[11pt]{article}
\usepackage{amsmath, amssymb, amsthm}
\usepackage[margin=1in]{geometry}

\newtheorem{definition}{Definition}
\newtheorem{proposition}{Proposition}
\newtheorem{theorem}{Theorem}
\newtheorem{corollary}{Corollary}

\title{Correctness of Tiled Union-Find for the Gaussian Moat Problem}
\author{}
\date{}

\begin{document}
\maketitle

%% ============================================================
%% SECTION 1 — Framework
%% ============================================================
\section{Framework}

\subsection{Acting Set}

Let $V$ be a finite set. Let $E$ be a set of unordered pairs from $V$
--- each element of $E$ is a pair $\{u, v\}$ with $u, v \in V$,
$u \ne v$. We call $(V, E)$ a graph.

For the Gaussian moat problem, $V$ is the set of Gaussian primes in a
bounded region of $\mathbb{Z}[i]$, and $\{p, q\} \in E$ iff
$\|p - q\|^2 \le K$ for a fixed threshold $K$.
The results below hold for any finite graph.

\subsection{Union-Find as Closure Operator}

Define $\operatorname{cl}(E)$ as the equivalence relation on $V$
generated by $E$: the reflexive, symmetric, transitive closure of the
edge relation. That is, $u \sim v$ under $\operatorname{cl}(E)$ iff
there exists a path $u = w_0, w_1, \ldots, w_m = v$ with each
$\{w_j, w_{j+1}\} \in E$.

The equivalence classes of $\operatorname{cl}(E)$ are the
\textbf{connected components} of the graph.

\textbf{Union-Find} is an algorithm that computes $\operatorname{cl}(E)$
by processing edges incrementally. Its output depends only on $E$, not
on processing order.

\begin{proposition}
\label{prop:closure-operator}
$\operatorname{cl}$ is a closure operator on the lattice of symmetric
relations over $V$:
\begin{enumerate}
    \item \textbf{Extensive:} $E \subseteq \operatorname{cl}(E)$
          \quad--- every edge is a length-1 path.
    \item \textbf{Monotone:} $E_1 \subseteq E_2 \;\Rightarrow\;
          \operatorname{cl}(E_1) \subseteq \operatorname{cl}(E_2)$
          \quad--- more edges, more paths.
    \item \textbf{Idempotent:}
          $\operatorname{cl}(\operatorname{cl}(E)) = \operatorname{cl}(E)$
          \quad--- the transitive closure of an equivalence relation is
          itself.
\end{enumerate}
\end{proposition}

\subsection{Decomposition Theorem}

\begin{theorem}[Closure Decomposition]
\label{thm:decomposition}
Let $\operatorname{cl}$ be a closure operator.
For any $A_1, \ldots, A_n$:
\[
    \operatorname{cl}(A_1 \cup \cdots \cup A_n)
    \;=\;
    \operatorname{cl}\!\bigl(
        \operatorname{cl}(A_1) \cup \cdots \cup \operatorname{cl}(A_n)
    \bigr).
\]
\end{theorem}

\begin{proof}
$(\subseteq)$\;
$A_i \subseteq \operatorname{cl}(A_i)$ by extensivity, so
$\bigcup A_i \subseteq \bigcup \operatorname{cl}(A_i)$, so
$\operatorname{cl}(\bigcup A_i) \subseteq
\operatorname{cl}(\bigcup \operatorname{cl}(A_i))$ by monotonicity.

\smallskip\noindent
$(\supseteq)$\;
$A_i \subseteq \bigcup A_i$ gives
$\operatorname{cl}(A_i) \subseteq \operatorname{cl}(\bigcup A_i)$
by monotonicity, so
$\bigcup \operatorname{cl}(A_i) \subseteq \operatorname{cl}(\bigcup A_i)$,
so $\operatorname{cl}(\bigcup \operatorname{cl}(A_i)) \subseteq
\operatorname{cl}(\operatorname{cl}(\bigcup A_i))
= \operatorname{cl}(\bigcup A_i)$
by monotonicity then idempotence.
\end{proof}

\subsection{Application to Tiled Union-Find}

Let $\{V_1, \ldots, V_n\}$ be subsets of $V$, each inducing an edge set
$E_i = \bigl\{\{u,v\} \in E : u, v \in V_i\bigr\}$. If

\[
    E_1 \cup \cdots \cup E_n = E
    \qquad\text{(\textbf{edge completeness})},
\]

then Theorem~\ref{thm:decomposition} gives directly:

\[
    \operatorname{cl}(E)
    \;=\;
    \operatorname{cl}\!\bigl(
        \operatorname{cl}(E_1) \cup \cdots \cup \operatorname{cl}(E_n)
    \bigr).
\]

Compute connected components within each subset independently, then
merge with one final closure step. The result equals the global
connected components.

In the tiling pipeline, each tile runs UF over its local vertices,
producing $\operatorname{cl}(E_i)$. The compositor initializes a fresh
UF over $V$ and, for each tile, unions all vertices sharing a local
component --- implementing the outer $\operatorname{cl}$ of the theorem.

\medskip\noindent
Two claims complete the argument:

\begin{enumerate}
    \item \textbf{Edge completeness} (\S\ref{sec:edge-completeness}):
          collar width $C \ge \lceil\sqrt{K}\,\rceil$ guarantees
          $\bigcup E_i = E$.
    \item \textbf{Port fidelity} (\S\ref{sec:port-fidelity}):
          the port/group encoding faithfully represents
          $\operatorname{cl}(E_i)$ at tile boundaries, so the
          compositor's merge step correctly implements the outer
          $\operatorname{cl}$.
\end{enumerate}

%% ============================================================
%% SECTION 2 — Edge Completeness (next)
%% ============================================================
\section{Edge Completeness}
\label{sec:edge-completeness}

% Claim: with collar C >= ceil(sqrt(K)), every edge {p,q} in E
% has both endpoints in at least one tile's vertex set V_i.
% Proof: max single-coordinate offset between connected primes
% is <= floor(sqrt(K)) < C.

%% ============================================================
%% SECTION 3 — Port Fidelity (next)
%% ============================================================
\section{Port Fidelity}
\label{sec:port-fidelity}

% Claim: the 128-byte TileOp (group labels + boundary ports)
% encodes exactly the boundary-to-boundary reachability within
% cl(E_i) — the information the outer cl needs.

\end{document}
